\chapter{Points, Distances and Families of Sets}\label{conjectures:ch:geometric}

\begin{center}\itshape
Colour the points of the plane so that no two points at distance one get the same colour. How few colours will do?\\
\upshape(Edward Nelson's question of 1950, in our words)
\end{center}

In 1950 Edward Nelson, a student at the University of Chicago, asked a question that a child can understand. Colour every point of the plane so that
two points at distance exactly $1$ never have the same colour. How many colours do you need? The same year John Isbell showed that seven colours are
enough. Four are needed, as Leo and William Moser's seven-point graph shows in three lines. And there the problem stood for sixty-eight years, with the
answer somewhere in $\{4,5,6,7\}$. In 2018 Aubrey de Grey, a biologist best known for his research on ageing, found a unit-distance graph that cannot be
coloured with four colours; his first graph had $20\,425$ vertices, the published one $1581$, and the checking was done by computer. The answer is now
known to be $5$, $6$ or $7$, and nobody has a guess that the experts agree on.

This chapter is about four problems of this kind, where the statement is about points or sets and fits in a line. The \emph{Hadwiger--Nelson problem}
asks for the chromatic number of the plane. The \emph{unit distance problem} of Erd\H os asks how many pairs among $n$ points of the plane can be at
distance $1$. The \emph{happy ending problem} asks how many points in general position force a convex $n$-gon. \emph{Frankl's conjecture} says that in a
family of sets closed under unions, some element lies in at least half of the sets. For each one we prove the classical bounds in full, and for the last
one we prove, with a short entropy argument, that some element lies in more than $15\%$ of the sets: a result of the kind that nobody could prove before
November 2022.

\begin{think}
Colour the plane with two colours. Can you avoid two points at distance $1$ with the same colour? Look at the three corners of an equilateral triangle
of side $1$. Now try three colours: where would you start looking for a contradiction?
\end{think}

\section{The chromatic number of the plane}

The \emph{unit-distance graph} of the plane has a vertex for every point of $\mathbb R^2$ and an edge between two points at distance $1$. Its chromatic
number $\chi(\mathbb R^2)$ is the least number of colours in a colouring of the plane with no monochromatic edge. A finite set of points with the unit
distances among them is a \emph{unit-distance graph}; if one of them needs $k$ colours, so does the plane.

\begin{theorem}[Moser and Moser, 1961; Isbell, 1950]\label{conjectures:thm:hn}
$4\le\chi(\mathbb R^2)\le7$.
\end{theorem}
\begin{proof}
\emph{Lower bound.} Let $A=(0,0)$. Let $B$ and $C$ be the two points at distance $1$ from $A$ and from each other, and $D$ the point at distance $1$ from
$B$ and $C$ other than $A$: the rhombus $ABDC$ is made of two equilateral triangles, and $|AD|=\sqrt3$. Rotate the rhombus about $A$ by the angle
$\theta$ with $2\sqrt3\sin\frac\theta2=1$, that is $\theta=2\arcsin\frac1{2\sqrt3}\approx33.56^\circ$, to get a second rhombus $AEGF$; then $|DG|=1$,
since $D$ and $G$ are on the circle of radius $\sqrt3$ about $A$, at angle $\theta$ apart. The seven points and eleven unit distances
$AB$, $AC$, $BC$, $BD$, $CD$, $AE$, $AF$, $EF$, $EG$, $FG$, $DG$ form the \emph{Moser spindle} (Figure~\ref{conjectures:fig:spindle}). Suppose it is coloured with
three colours. $A$, $B$, $C$ form a triangle, so they get the three colours; $D$ is joined to $B$ and $C$, so it gets the colour of $A$. In the same way
$G$ gets the colour of $A$. But $D$ and $G$ are joined: a contradiction. So $\chi(\mathbb R^2)\ge4$.

\emph{Upper bound.} Tile the plane by regular hexagons of side $s$, with $0.39<s<\frac12$ (say $s=0.45$). The centres form a triangular lattice
$\{ia_1+ja_2:i,j\in\mathbb Z\}$ with $|a_1|=|a_2|=\sqrt3\,s$ and an angle of $60^\circ$ between $a_1$ and $a_2$. Give the hexagon centred at $ia_1+ja_2$
the colour $i+3j\bmod7$; the boundary points go to any one of the hexagons that contain them. Two points in the same hexagon are at distance at most its
diameter $2s<1$. Two hexagons of the same colour have centres differing by $ia_1+ja_2\ne0$ with $7\mid i+3j$, and
\[
|ia_1+ja_2|^2=3s^2(i^2+ij+j^2).
\]
The nonzero values of $i^2+ij+j^2$ below $7$ are $1$, $3$ and $4$ (Exercise~\ref{conjectures:ex:hexcolour}), taken at eighteen vectors:
\[
\begin{aligned}
1&:\ \pm(1,0),\ \pm(0,1),\ \pm(1,-1),\\
3&:\ \pm(1,1),\ \pm(2,-1),\ \pm(1,-2),\\
4&:\ \pm(2,0),\ \pm(0,2),\ \pm(2,-2).
\end{aligned}
\]
The values of $i+3j$ modulo $7$ at these vectors are $\pm1,\pm3,\pm2$, then $\pm4,\pm1,\pm2$, then $\pm2,\pm6,\pm4$: never $0$.
So same-coloured centres are at distance at least $\sqrt{21}\,s$, and every point of a hexagon is within $s$ of its centre, so two points of the same
colour in different hexagons are at distance at least $(\sqrt{21}-2)s>1$, because $s>0.39>\frac1{\sqrt{21}-2}=0.3871\ldots$ No two points at distance $1$
have the same colour.
\end{proof}

\begin{figure}[ht]
\centering
\fig{conjectures-ch08-spindle}
\caption{Left: the Moser spindle, seven points with eleven unit distances (the edges). In a $3$-colouring, $D$ and $G$ would both get the colour of $A$,
but $DG$ (red) has length $1$. The colours shown are a $3$-colouring of everything except the edge $DG$. Right: Isbell's $7$-colouring by hexagons of
side $0.45$; the hexagon with lattice coordinates $(i,j)$ gets colour $i+3j\bmod7$, so each hexagon and its six neighbours have seven different colours.}
\label{conjectures:fig:spindle}
\end{figure}

For a long time the two bounds were all there was. The problem is named after Hugo Hadwiger, who had studied closely related covering questions, and
Nelson; it reached a wide public through Martin Gardner's column in \emph{Scientific American} in 1960.

\begin{remark}[five colours, and what remains]\label{conjectures:rem:degrey}
In April 2018 de Grey posted a proof that $\chi(\mathbb R^2)\ge5$: a unit-distance graph that cannot be $4$-coloured, built from many copies of the Moser
spindle and of a hexagonal piece of the triangular lattice, with the non-colourability checked by a computer (a SAT solver). The collaborative
\emph{Polymath16} project, hosted on Dustin Mixon's blog, shrank the graph; the smallest known $5$-chromatic unit-distance
graph, found by Jaan Parts in 2020, has $509$ vertices. Two other facts frame the problem. By the de Bruijn--Erd\H os theorem (which uses the axiom of
choice), $\chi(\mathbb R^2)$ is the largest chromatic number of a \emph{finite} unit-distance graph, so the question is really about finite
configurations. And if the colour classes are required to be measurable, at least five colours are needed, by a theorem of Kenneth Falconer (1981). No
one knows a $6$-chromatic unit-distance graph, and no one knows a $6$-colouring of the plane with no monochromatic unit distance.
\end{remark}

\paragraph{Why it is hard.} The lower bounds need finite graphs, and the only way known to find them is to glue many copies of rigid small graphs and
let a computer check that the result cannot be coloured; each extra colour seems to need a graph larger by orders of magnitude. The upper bound needs a
colouring of the whole plane, and the hexagon colouring of 1950 has resisted every attempt at improvement: a $6$-colouring would have to break the
pleasant symmetry of all known good colourings.

\section{How many unit distances?}

Let $u(n)$ be the largest number of pairs at distance exactly $1$ among $n$ points of the plane. In 1946 Paul Erd\H os observed that a suitable piece of
the integer lattice, rescaled, has $n^{1+c/\log\log n}$ unit distances, and asked whether this is the truth. His \emph{unit distance conjecture} says
that $u(n)\le n^{1+\varepsilon}$ for every $\varepsilon>0$ and $n$ large. A simple construction already beats any linear function
(Exercise~\ref{conjectures:ex:unitcube}). For the upper bound, Spencer, Szemer\'edi and Trotter proved $u(n)=O(n^{4/3})$ in 1984, and in 1997 L\'aszl\'o Sz\'ekely
found the proof below. Nothing better is known.

We use the crossing lemma (Volume~3, Exercise~12.4): a simple graph with $v$ vertices and $m\ge4v$ edges, drawn in the plane with edges as curves, has at
least $\frac{m^3}{64v^2}$ pairs of crossing edges.

\begin{theorem}[Spencer, Szemer\'edi and Trotter, 1984; Sz\'ekely's proof]\label{conjectures:thm:unitdist}
For $n\ge1$, $u(n)<4n^{4/3}+n$.
\end{theorem}
\begin{proof}
Let $P$ be a set of $n$ points and $u$ the number of unit distances among them. For $p\in P$ let $C_p$ be the circle of radius $1$ about $p$ and
$k_p=|C_p\cap P|$. Each unit distance $\{p,q\}$ puts $q$ on $C_p$ and $p$ on $C_q$, so $\sum_pk_p=2u$.

Draw a graph $G$ with vertex set $P$: on each circle $C_p$ with $k_p\ge3$, join consecutive points of $P$ along the circle by the arc between them. This
gives $k_p$ arcs on $C_p$, none of which passes through a point of $P$. The circles with $k_p\le2$ are dropped, and they account for at most $2n$ of the
sum $\sum k_p$, so there are at least $2u-2n$ arcs. Two distinct points lie on at most two unit circles (the centres are the at most two points at
distance $1$ from both), so each pair of points is joined by at most two arcs; keeping one arc for each pair gives a simple graph with
$m\ge u-n$ edges. Two arcs can only cross at a point where their circles meet; two circles meet in at most two points, and at such a point at most one
arc of each circle passes. So the drawing has at most $2\binom n2<n^2$ crossing pairs.

If $m\ge4n$, the crossing lemma (with $v\le n$ vertices) gives $\frac{m^3}{64n^2}<n^2$, so $m<4n^{4/3}$. If $m<4n$, then also $m<4n^{4/3}$. In both cases
$u\le m+n<4n^{4/3}+n$.
\end{proof}

The same method gives the Szemer\'edi--Trotter theorem on incidences between points and lines; the problem book P1 (Problem~8.3) develops both. The
exponent $\frac43$ is the truth for lines, by a grid example, which is why the method cannot do better for circles without a new idea.

\paragraph{Why it is hard.} The proof uses only two facts about unit circles: two of them meet in at most two points, and two points lie on at most two of
them. These facts are also true for curves where the bound $n^{4/3}$ is attained (for instance, for suitable families of parabolas or for points and
lines), so any improvement must use something specific to Euclidean circles, such as the arithmetic of the distances. In other norms, and in dimension
$3$, the answer changes; in dimension $4$ it becomes quadratic (Exercise~\ref{conjectures:ex:lenz}).

\section{The happy ending problem}

Around 1933, in Budapest, Esther Klein observed that among five points of the plane, no three on a line, four are always the vertices of a
convex quadrilateral (Volume~3, Exercise~A.13). She asked whether, for every $n$, enough points force a convex $n$-gon. George Szekeres and
Paul Erd\H os proved it in 1935, and Erd\H os called it the \emph{happy ending problem} because Klein and Szekeres married in 1937 (Volume~3, Appendix~B, tells the story). Let $f(n)$ be the least
number such that any $f(n)$ points in general position (no three on a line) contain $n$ points in \emph{convex position}, the vertices of a convex
$n$-gon. Volume~3, Appendix~B proves that $f(n)$ is finite with Ramsey's theorem; here is the much better bound of the original paper, by a cleverer argument.

Rotate the plane so that no two points have the same first coordinate, and order the points from left to right. A \emph{$k$-cup} is a sequence of $k$
points $p_1,\dots,p_k$ from left to right such that the slopes of $p_1p_2,p_2p_3,\dots,p_{k-1}p_k$ increase; a \emph{$k$-cap} is the same with the slopes
decreasing. Three points from left to right form a $3$-cup or a $3$-cap, since they are not on a line.

\begin{lemma}\label{conjectures:lem:cupconvex}
The points of a cup, or of a cap, are in convex position.
\end{lemma}
\begin{proof}
Let $p_1,\dots,p_k$ be a cup and $\varphi$ the piecewise linear function on $[x_1,x_k]$ whose graph is the path $p_1p_2\cdots p_k$. Its slopes increase,
so it is convex, and the region $R$ between its graph and the segment $p_1p_k$ is convex. Each $p_i$ is an extreme point of $R$: for $1<i<k$, the slope of
$\varphi$ jumps strictly at $x_i$, so $p_i$ is not the midpoint of two points of $R$; and $p_1$, $p_k$ are the points of $R$ with the smallest and
largest first coordinate. An extreme point of a convex set is not in the convex hull of other points of the set, so no $p_i$ is in the convex hull of
the others. For a cap, reflect in the horizontal axis.
\end{proof}

\begin{theorem}[Erd\H os and Szekeres, 1935]\label{conjectures:thm:cupcap}
Let $g(k,l)$ be the least $N$ such that any $N$ points in general position with distinct first coordinates contain a $k$-cup or an $l$-cap. Then for
$k,l\ge2$,
\[
g(k,l)\le\binom{k+l-4}{k-2}+1,
\]
and consequently $f(n)\le\binom{2n-4}{n-2}+1$.
\end{theorem}
\begin{proof}
Any two points form a $2$-cup and a $2$-cap, so $g(2,l)=g(k,2)=2$, as the formula says. We show $g(k,l)\le g(k-1,l)+g(k,l-1)-1$ for $k,l\ge3$;
Pascal's rule then gives the bound by induction on $k+l$.

Let $P$ have $g(k-1,l)+g(k,l-1)-1$ points, and let $E$ be the set of points of $P$ that are the right end of some $(k-1)$-cup. If $P\setminus E$ has at
least $g(k-1,l)$ points, it contains an $l$-cap, and we are done, or a $(k-1)$-cup, whose right end is in $E$: impossible. Otherwise $E$ has at least
$g(k,l-1)$ points, so it contains a $k$-cup, and we are done, or an $(l-1)$-cap $q=q_1,q_2,\dots,q_{l-1}$. Its left end $q$ is in $E$, the right end of a
$(k-1)$-cup ending $p,q$. The points $p,q,q_2$ are not on a line. If the slope of $qq_2$ is larger than that of $pq$, the cup followed by $q_2$ is a
$k$-cup. Otherwise it is smaller, and $p$ followed by the cap is an $l$-cap.

Finally, $n$ points forming an $n$-cup or an $n$-cap are in convex position (Lemma~\ref{conjectures:lem:cupconvex}), so $f(n)\le g(n,n)$.
\end{proof}

For $n=5$ the bound gives $f(5)\le\binom63+1=21$, far above the truth. The bound $g(k,l)$ itself is exact (Erd\H os and Szekeres gave the matching
example; Exercise~\ref{conjectures:ex:cupcap} does the first cases), so the loss is in passing from cups and caps to general convex polygons.

\begin{proposition}\label{conjectures:prop:eight}
There are eight points in general position with no five in convex position. Hence $f(5)\ge9$.
\end{proposition}
\begin{proof}
Let $\rho$ be the rotation by $90^\circ$ about the origin, $P_i=\rho^i(6,0)$ and $Q_i=\rho^i(1,2)$ for $i=0,1,2,3$ (indices modulo $4$); see
Figure~\ref{conjectures:fig:eight}. No three of the points are on a line (a direct check of the $56$ triples, done numerically). The rotation $\rho$ permutes
the points and preserves convex position, so it is enough to check one five-point set in each orbit. A five-point set is not in convex position as
soon as one of its points lies inside the triangle of three others, and we exhibit such a point in each case.

The triangle of three of the $P_i$ is the half of the square $P_0P_1P_2P_3$ cut off by a diagonal. $Q_0=(1,2)$ lies in the open quadrant $x,y>0$, so it is
inside the two half-squares $P_3P_0P_1$ ($x\ge0$) and $P_0P_1P_2$ ($y\ge0$); by rotation, $Q_i$ is inside the triangle of three $P$'s unless the
missing $P$ is $P_i$ or $P_{i+1}$.

\emph{Four $P$'s.} Any $Q_i$ is inside one of the half-squares.

\emph{Three $P$'s.} Up to rotation the missing one is $P_3$, and the two $Q$'s must both lie outside the half-square $P_0P_1P_2$, so they are $Q_2$ and
$Q_3$. In $\{P_0,P_1,P_2,Q_2,Q_3\}$, the point $Q_3=(2,-1)$ lies inside the triangle $P_0P_1Q_2$, with $Q_2=(-1,-2)$.

\emph{Two $P$'s, adjacent.} Up to rotation they are $P_0,P_1$, with three of the four $Q$'s. The point $Q_0$ lies inside the triangles $P_0P_1Q_1$ and
$P_0P_1Q_2$, and $Q_3$ inside $P_0P_1Q_2$. A triple of $Q$'s contains $Q_0$ and $Q_1$, or $Q_0$ and $Q_2$, or is $\{Q_1,Q_2,Q_3\}$; one of the three facts
applies.

\emph{Two $P$'s, opposite.} Up to rotation they are $P_0,P_2$. The point $Q_1$ lies inside $P_0P_2Q_0$, which settles the triples containing $Q_0$ and $Q_1$;
the half-turn $\rho^2$ fixes $\{P_0,P_2\}$ and maps these triples to the other two, $\{Q_2,Q_3,Q_0\}$ and $\{Q_2,Q_3,Q_1\}$.

\emph{One $P$.} Up to rotation it is $P_0$, with all four $Q$'s, and $Q_3$ lies inside $P_0Q_0Q_2$.

A five-point set has at most four $Q$'s, so these cases are all. Each ``lies inside'' is a check of three signs of determinants
(Exercise~\ref{conjectures:ex:pentwitness}).
\end{proof}

\begin{figure}[ht]
\centering
\fig{conjectures-ch08-eight}
\caption{Eight points with no convex pentagon: the square $P_0P_1P_2P_3$ with $P_0=(6,0)$, and a smaller tilted square $Q_0Q_1Q_2Q_3$ with
$Q_0=(1,2)$, both invariant under the rotation by $90^\circ$. The tilt matters: with $Q_0$ on a diagonal of the big square, three points would be on a
line.}
\label{conjectures:fig:eight}
\end{figure}

\begin{remark}[what is known about $f(n)$]\label{conjectures:rem:esz}
$f(3)=3$, $f(4)=5$ (Klein), $f(5)=9$, a result credited to Endre Makai, and $f(6)=17$, proved by George Szekeres and Lindsay Peters in 2006 by a computer
search. In 1961 Erd\H os and Szekeres constructed, for every $n$, a set of $2^{n-2}$ points in general position with no convex $n$-gon, by gluing
cup-cap examples, and conjectured that $f(n)=2^{n-2}+1$. In 2017 Andrew Suk proved $f(n)\le2^{n+O(n^{2/3}\log n)}$, so that $f(n)=2^{n+o(n)}$; the
exact conjecture is open from $n=7$ on.
\end{remark}

\paragraph{Why it is hard.} Cups and caps reduce the problem to sequences, where the counting is exact, but convex polygons can mix the two shapes in
many ways. Suk's proof needs a delicate decomposition of the point set and the Erd\H os--Szekeres bound on cups and caps in each piece, with an error
term too large to see the exact value. For $n=7$ a direct computer search, as for $n=6$, is far out of reach: the configurations of $32$ points are too
many, even up to the right equivalence.

\section{Union-closed families}

A family $\mathcal F$ of subsets of a finite set is \emph{union-closed} if $A\cup B\in\mathcal F$ whenever $A,B\in\mathcal F$. The \emph{frequency} of an
element $x$ is the proportion of sets of $\mathcal F$ that contain $x$.

\begin{center}
\emph{Frankl's conjecture.} If $\mathcal F$ is a finite union-closed family containing at least one nonempty set, some element has frequency at least
$\frac12$.
\end{center}

The power set of $\{1,\dots,n\}$ shows that $\frac12$ cannot be improved. As Henning Bruhn and Oliver Schaudt tell the story in their survey, P\'eter
Frankl formulated the conjecture in 1979, on a journey from Paris to Montreal, and mentioned it to Ronald Graham, who spread it; some authors consider it
folklore from earlier in the 1970s. It has been checked for families on a ground set of at most $12$ elements and for families of at most $50$ sets, and
proved for many special classes. Until 2022, nobody could prove that some element has frequency at least $c$ for any fixed $c>0$: the best general
bounds tended to $0$ as the family grew.

\subsection{Two easy cases}

\begin{proposition}\label{conjectures:prop:ucsmall}
Let $\mathcal F$ be a finite union-closed family.
\begin{enumerate}[label=(\alph*)]
\item If $\{x\}\in\mathcal F$, then $x$ has frequency at least $\frac12$.
\item If $\{x,y\}\in\mathcal F$, then $x$ or $y$ has frequency at least $\frac12$.
\end{enumerate}
\end{proposition}
\begin{proof}
(a) The map $S\mapsto S\cup\{x\}$ sends the sets of $\mathcal F$ not containing $x$ into the sets of $\mathcal F$ containing $x$ (union-closed), and it is
injective on them, since $S=(S\cup\{x\})\setminus\{x\}$. So at least half of the sets contain $x$.

(b) Split $\mathcal F$ into $\mathcal F_\emptyset$, $\mathcal F_x$, $\mathcal F_y$, $\mathcal F_{xy}$ according to which of $x,y$ a set contains. As in (a),
$S\mapsto S\cup\{x,y\}$ is injective from $\mathcal F_\emptyset$ to $\mathcal F_{xy}$, so $|\mathcal F_\emptyset|\le|\mathcal F_{xy}|$. The numbers of sets
containing $x$ and containing $y$ add up to
\[
|\mathcal F_x|+|\mathcal F_y|+2|\mathcal F_{xy}|\ge|\mathcal F_x|+|\mathcal F_y|+|\mathcal F_{xy}|+|\mathcal F_\emptyset|=|\mathcal F|,
\]
so one of them is at least $\frac{|\mathcal F|}2$.
\end{proof}

For a $3$-set $\{x,y,z\}$ the same counting no longer closes up: the injection into the sets containing all three is not enough to beat the
sets that contain exactly one of them. The two cases show the shape of most partial results: some small structure in $\mathcal F$ is pushed around by unions.

\subsection{Entropy}

In November 2022 Justin Gilmer found a completely different approach, through entropy. In this section logarithms are natural, so entropies are in
nats; Volume~3 used base $2$, and changing the base multiplies every entropy by the same constant, which changes nothing below. For a random variable $X$
with finitely many values, $H(X)=\sum_xP(X=x)\ln\frac1{P(X=x)}$, and $h(p)=p\ln\frac1p+(1-p)\ln\frac1{1-p}$ is the entropy of a bit that is $0$ with
probability $p$. By Gibbs' inequality (Volume~3, Proposition~4.16), $H(X)\le\ln m$ if $X$ takes at most $m$ values, with equality for the uniform law.

For random variables $X,Y$ with finitely many values, let $H(X\mid Y=y)$ be the entropy of the law of $X$ given $Y=y$, and
$H(X\mid Y)=\sum_yP(Y=y)H(X\mid Y=y)$.

\begin{lemma}\label{conjectures:lem:condent}
\begin{enumerate}[label=(\alph*)]
\item (Chain rule.) $H(X,Y)=H(Y)+H(X\mid Y)$, and $H(X_1,\dots,X_n)=\sum_{i=1}^nH(X_i\mid X_1,\dots,X_{i-1})$.
\item (Conditioning reduces entropy.) $H(X\mid Y,W)\le H(X\mid Y)$. If $Y$ is a function of $W$, then $H(X\mid W)\le H(X\mid Y)$.
\end{enumerate}
\end{lemma}
\begin{proof}
(a) Write $P(X=x,Y=y)=P(Y=y)P(X=x\mid Y=y)$, take $\ln\frac1{\,\cdot\,}$, which turns the product into a sum, multiply by $P(X=x,Y=y)$ and add over
$(x,y)$. The second formula follows by induction, applied to $(X_1,\dots,X_{n-1})$ and $X_n$.

(b) Fix $y$, and let $\lambda_w=P(W=w\mid Y=y)$ and $q_w(x)=P(X=x\mid Y=y,W=w)$. Then $P(X=x\mid Y=y)=\sum_w\lambda_wq_w(x)$, and since
$\eta(t)=t\ln\frac1t$ is concave on $[0,1]$ (its second derivative is $-\frac1t$), Jensen's inequality gives
$\eta\bigl(\sum_w\lambda_wq_w(x)\bigr)\ge\sum_w\lambda_w\eta(q_w(x))$. Adding over $x$: $H(X\mid Y=y)\ge\sum_w\lambda_wH(X\mid Y=y,W=w)$. Multiply by
$P(Y=y)$ and add over $y$. If $Y$ is a function of $W$, then conditioning on $(Y,W)$ is conditioning on $W$.
\end{proof}

Gilmer's idea is to take two independent uniform random sets $A,B$ in $\mathcal F$. Their union is again in $\mathcal F$, so it cannot have more entropy
than $A$. But if all elements are rare, a union of two independent sets is more random than one set, because each element gets two chances. The whole
proof is the comparison of these two entropies, coordinate by coordinate.

\begin{theorem}[Gilmer's method]\label{conjectures:thm:gilmer}
Suppose that $c>1$ is such that $\mathbb E\,h(UV)\ge c\,\mathbb E[U]\,\mathbb E\,h(U)$ for every pair of independent, identically distributed random
variables $U,V$ with finitely many values in $[0,1]$. Then in every finite union-closed family $\mathcal F$ with at least two sets, some element has
frequency at least $1-\frac1c$.
\end{theorem}
\begin{proof}
Let the ground set be $\{1,\dots,n\}$ and identify a set with its indicator vector in $\{0,1\}^n$. Let $A,B$ be independent and uniform on $\mathcal F$,
and $C=A\cup B$, with coordinates $C_i=\max(A_i,B_i)$. Write $A_{<i}=(A_1,\dots,A_{i-1})$, and similarly for $B$ and $C$. Suppose that every element has
frequency at most $p$, and let us show that $c(1-p)\le1$.

For a vector $a\in\{0,1\}^{i-1}$ with $P(A_{<i}=a)>0$, let $u_i(a)=P(A_i=0\mid A_{<i}=a)$, and put $U_i=u_i(A_{<i})$ and $V_i=u_i(B_{<i})$. These are
independent with the same law, and $\mathbb E[U_i]=P(A_i=0)\ge1-p$. By the chain rule,
\[
\ln|\mathcal F|=H(A)=\sum_{i=1}^nH(A_i\mid A_{<i})=\sum_{i=1}^n\mathbb E\,h(U_i).
\]
For $C$, the vector $C_{<i}$ is a function of $(A_{<i},B_{<i})$, so by Lemma~\ref{conjectures:lem:condent}(b),
$H(C_i\mid C_{<i})\ge H(C_i\mid A_{<i},B_{<i})$. Given $A_{<i}=a$ and $B_{<i}=b$, the bits $A_i$ and $B_i$ are independent (because $A$ and $B$ are), and
$C_i=0$ if and only if both are $0$, which has probability $u_i(a)u_i(b)$. So $H(C_i\mid A_{<i},B_{<i})=\mathbb E\,h(U_iV_i)$, and by the hypothesis,
\[
H(C)=\sum_{i=1}^nH(C_i\mid C_{<i})\ge\sum_{i=1}^n\mathbb E\,h(U_iV_i)\ge c\sum_{i=1}^n\mathbb E[U_i]\,\mathbb E\,h(U_i)\ge c(1-p)\ln|\mathcal F|.
\]
But $C$ takes values in $\mathcal F$ because $\mathcal F$ is union-closed, so $H(C)\le\ln|\mathcal F|$ by Gibbs' inequality. As $|\mathcal F|\ge2$,
$\ln|\mathcal F|>0$, and $c(1-p)\le1$, that is $p\ge1-\frac1c$. Taking for $p$ the largest frequency proves the theorem.
\end{proof}

Everything now rests on an inequality between numbers. The weak form $h(xy)\ge x\,h(y)$ is just concavity: $h(xy)=h\bigl(x\cdot y+(1-x)\cdot0\bigr)\ge
x\,h(y)+(1-x)h(0)$. It gives the hypothesis with $c=1$ only, which proves nothing. We need to gain a constant factor over it, and here is a version with a
short proof.

\begin{lemma}\label{conjectures:lem:hxy}
For $x,y\in[0,1]$,
\[
h(xy)\ge\frac1{1+\ln2}\bigl(x\,h(y)+y\,h(x)\bigr).
\]
\end{lemma}
\begin{proof}
If $x$ or $y$ is $0$, both sides are $0$. For $0<t\le1$ put $t=e^{-a}$ with $a\ge0$; then
\[
\frac{h(t)}t=\ln\frac1t+\frac{1-t}t\ln\frac1{1-t}=a+K(a),\qquad K(a)=(e^a-1)\ln\frac1{1-e^{-a}}\quad(K(0)=0).
\]
\emph{A series for $K$.} With $\ln\frac1{1-z}=\sum_{m\ge1}\frac{z^m}m$ for $z=e^{-a}<1$,
\[
K(a)=\sum_{m\ge1}\frac{e^{-(m-1)a}-e^{-ma}}m=1-\sum_{m\ge1}\frac{e^{-ma}}{m(m+1)},
\]
by regrouping the terms of each power of $e^{-a}$ (the series converge absolutely), and this also holds at $a=0$, as $\sum\frac1{m(m+1)}=1$. In
particular $K\ge0$, and for $a,b\ge0$ the \emph{defect} is
\[
D(a,b)=K(a)+K(b)-K(a+b)=\sum_{m\ge1}\frac{(1-e^{-ma})(1-e^{-mb})}{m(m+1)}.
\]
\emph{The defect is at most $(a+b)\ln2$.} The function $\ell(z)=\ln(1-e^{-z})$ is concave on $]0,\infty[$, since $\ell'(z)=\frac1{e^z-1}$ decreases. So
$(1-e^{-ma})(1-e^{-mb})\le(1-e^{-ms})^2$ with $s=\frac{a+b}2$, and $D(a,b)\le\sum_m\frac{(1-e^{-ms})^2}{m(m+1)}=:F(s)$. We show $F(s)\le2s\ln2$. Both sides
vanish at $s=0$; differentiating term by term (legitimate for $s>0$, where the differentiated series converges uniformly near $s$), with $t=e^{-s}$,
\[
F'(s)=2\sum_{m\ge1}\frac{t^m-t^{2m}}{m+1}=2\Bigl(\frac1t\ln\frac1{1-t}-\frac1{t^2}\ln\frac1{1-t^2}\Bigr)=\frac2{t^2}\Bigl(\ln(1+t)+(1-t)\ln(1-t)\Bigr),
\]
using $\sum_{m\ge1}\frac{z^m}{m+1}=\frac1z\ln\frac1{1-z}-1$. Let $Q(t)=t^2\ln2-\ln(1+t)-(1-t)\ln(1-t)$ on $[0,1]$. Then $Q(0)=Q(1)=0$,
$Q'(t)=2t\ln2+\frac t{1+t}+\ln(1-t)$, $Q'(0)=0$, $Q''(t)=2\ln2+\frac1{(1+t)^2}-\frac1{1-t}$, and $Q'''(t)=-\frac2{(1+t)^3}-\frac1{(1-t)^2}<0$. So $Q''$
decreases from $Q''(0)=2\ln2>0$ to $-\infty$: $Q'$ first increases from $0$, then decreases to $-\infty$, so it is positive and then negative; $Q$
increases and then decreases, and since $Q(0)=Q(1)=0$, $Q\ge0$ on $[0,1]$. Hence $F'(s)\le2\ln2$, and $F(s)\le2s\ln2$.

\emph{Conclusion.} Let $x=e^{-a}$, $y=e^{-b}$, and $c'=\frac1{1+\ln2}$, so that $1-c'=c'\ln2$. Dividing the claim by $xy$, we must show
$a+b+K(a+b)\ge c'\bigl(a+K(a)+b+K(b)\bigr)$. The difference of the two sides is
\[
(1-c')(a+b)+(1-c')K(a+b)-c'D(a,b)\ge(1-c')(a+b)-c'(a+b)\ln2=0,
\]
because $K\ge0$ and $D(a,b)\le(a+b)\ln2$.
\end{proof}

\begin{corollary}\label{conjectures:cor:ucconst}
In every finite union-closed family with at least two sets, some element has frequency at least $\frac{1-\ln2}2=0.1534\ldots$
\end{corollary}
\begin{proof}
For $U,V$ independent with the same law, Lemma~\ref{conjectures:lem:hxy} and independence give
\[
\begin{aligned}\mathbb E\,h(UV)&\ge\frac1{1+\ln2}\bigl(\mathbb E[U]\,\mathbb E\,h(V)+\mathbb E[V]\,\mathbb E\,h(U)\bigr)\\&=\frac2{1+\ln2}\,\mathbb E[U]\,\mathbb E\,h(U).\end{aligned}
\]
Apply
Theorem~\ref{conjectures:thm:gilmer} with $c=\frac2{1+\ln2}$: $1-\frac1c=1-\frac{1+\ln2}2=\frac{1-\ln2}2$.
\end{proof}

\begin{remark}[the race of November 2022, and the golden ratio]\label{conjectures:rem:ucrace}
Gilmer's paper proved the frequency bound $0.01$, with an inequality weaker than Lemma~\ref{conjectures:lem:hxy}. Within a week, several groups (Ryan Alweiss,
Brice Huang and Mark Sellke; Zachary Chase and Shachar Lovett; Will Sawin; Stijn Cambie) raised it to $\frac{3-\sqrt5}2=0.38196\ldots$ The key
inequality is
\[
h(xy)\ge\frac\varphi2\bigl(x\,h(y)+y\,h(x)\bigr),\qquad\varphi=\frac{1+\sqrt5}2,
\]
proved by Ravi Boppana in the 1980s for a question on Boolean formulas, with equality at $x=y=\frac1\varphi$; it gives Theorem~\ref{conjectures:thm:gilmer} with
$c=\varphi$ and $1-\frac1\varphi=\frac{3-\sqrt5}2$. Its proof is a page of careful calculus. The constant $\frac{3-\sqrt5}2$ is a natural barrier for
this exact form of the method (Exercise~\ref{conjectures:ex:ucgolden}); Sawin showed how to go slightly beyond it, and Jingbo Liu reached $0.38271$ in 2023, with a
proof that relies on numerically verified hypotheses. The conjectured $\frac12$ is still far away.
\end{remark}

\paragraph{Why it is hard.} The extremal families for Frankl's conjecture are power sets, where every element has frequency exactly $\frac12$, and many
other families come close; a proof must be exact on all of them. The entropy method loses because it compares one union $A\cup B$ with $A$, while the
conjecture needs to see the whole structure of $\mathcal F$. Before Gilmer, the methods were combinatorial (weights, averaging over special sets), and
none gave a constant at all.

\section{Reading a paper: Gilmer 2022}

Justin Gilmer's \emph{A constant lower bound for the union-closed sets conjecture} (arXiv:\allowbreak2211.09055) is short, and worth reading for its strategy rather
than its constant. The main theorem is proved by contradiction: assume every element has small frequency, and show that the union of two independent
copies of a random set from $\mathcal F$ has more entropy than the uniform law on $\mathcal F$ allows. The proof has the structure of
Theorem~\ref{conjectures:thm:gilmer}: the chain rule over the coordinates, conditioning on the past of both sets, and a one-coordinate inequality. The inequality is
where Gilmer's version is lossy, and where all the follow-up papers improved it. Read the paper with this chapter's Lemma~\ref{conjectures:lem:hxy} at hand, and
notice how much of the argument is independent of the constant: this is a sign that the method, not the estimate, was the discovery. Then read the
paper of Alweiss, Huang and Sellke, which explains clearly why $\frac{3-\sqrt5}2$ is where this form of the method stops.

\section{Open doors}

\begin{itemize}
\item Search, with a SAT solver, for $5$-chromatic unit-distance graphs inside the Moser lattice $\{a+b\omega_1+c\omega_3+d\omega_1\omega_3\}$ generated
by sixth roots of unity and their rotation by the spindle angle. How small can you make one?
\item Compute $u(n)$ exactly for $n\le10$ by searching through candidate configurations, and compare with the values in the literature.
\item Build Erd\H os and Szekeres's $2^{n-2}$-point sets without a convex $n$-gon for $n=5,6$ and check them by computer.
\item Run the entropy argument numerically: for random union-closed families, compute $H(A\cup B)$ and $H(A)$, and the coordinate-by-coordinate terms.
How close to equality is each step?
\end{itemize}

\section{Exercises}

\begin{exercise}[$\star$]\label{conjectures:ex:linetwo}
Show that the unit-distance graph of the real line has chromatic number $2$, and that the plane needs at least $3$ colours.
\end{exercise}

\begin{exercise}[$\star\star$]\label{conjectures:ex:spindle}
Show that the Moser spindle has no three pairwise non-adjacent vertices, and deduce again that it needs four colours. Show that removing any one of its
eleven edges leaves a $3$-colourable graph.
\end{exercise}

\begin{exercise}[$\star\star$]\label{conjectures:ex:hexcolour}
Show that $i^2+ij+j^2$ never takes the values $2$, $5$, $6$ for integers $i,j$, and find all $(i,j)$ where it takes the values $1$, $3$, $4$. In the
colouring of Theorem~\ref{conjectures:thm:hn}, check that the six neighbours of a hexagon and the hexagon itself have seven different colours. For which
sides $s$ does the proof work?
\end{exercise}

\begin{exercise}[$\star\star$]\label{conjectures:ex:unitcube}
Let $v_1,\dots,v_k$ be unit vectors of the plane such that the $2^k$ sums $\sum_{i\in S}v_i$, $S\subseteq\{1,\dots,k\}$, are distinct. Show that such
vectors exist, and that the $n=2^k$ sums determine at least $\frac n2\log_2n$ unit distances.
\end{exercise}

\begin{exercise}[$\star$]\label{conjectures:ex:lenz}
In $\mathbb R^4$, put $\frac n2$ points on the circle $\{(x,y,0,0):x^2+y^2=\frac12\}$ and $\frac n2$ on $\{(0,0,z,w):z^2+w^2=\frac12\}$ ($n$ even). Show
that they determine at least $\frac{n^2}4$ unit distances.
\end{exercise}

\begin{exercise}[$\star\star$]\label{conjectures:ex:cupcap}
Show that $g(3,l)=l$ for $l\ge2$, by showing that $l-1$ points on the graph of a concave function contain no $3$-cup and no $l$-cap. Deduce $g(k,3)=k$.
\end{exercise}

\begin{exercise}[$\star$]\label{conjectures:ex:pentwitness}
Check that $Q_3$ lies inside the triangle $P_0P_1Q_2$ and that $Q_1$ lies inside $P_0P_2Q_0$ in Proposition~\ref{conjectures:prop:eight}, using signs of
determinants.
\end{exercise}

\begin{exercise}[$\star$]\label{conjectures:ex:ucsharp}
(a) Show that in the power set of $\{1,\dots,n\}$ every element has frequency $\frac12$. (b) Find a union-closed family in which some element has
frequency below $\frac12$, and one which is not union-closed and in which every element has frequency $\frac13$. (c) Let $\mathcal F$ be union-closed with
$\emptyset\in\mathcal F$ and $N\ge3$ sets. Show that $\mathcal F\setminus\{\emptyset\}$ is union-closed, and that if Frankl's conjecture holds for
$\mathcal F$ it holds for $\mathcal F\setminus\{\emptyset\}$. In the other direction, what frequency in $\mathcal F$ do you get from frequency
$\frac12$ in $\mathcal F\setminus\{\emptyset\}$?
\end{exercise}

\begin{exercise}[$\star$]\label{conjectures:ex:ucdual}
Show that Frankl's conjecture is equivalent to: in every finite family $\mathcal G$ of at least two sets, closed under intersections, some element of
$\bigcup\mathcal G$ has frequency at most $\frac12$.
\end{exercise}

\begin{exercise}[$\star\star$]\label{conjectures:ex:ucgolden}
(a) Check that $h(x^2)=\varphi\,x\,h(x)$ at $x=\frac1\varphi$. (b) Show that the hypothesis of Theorem~\ref{conjectures:thm:gilmer} fails for every $c>\varphi$ (try
$U$ and $V$ constant). (c) Why does this not show that Frankl's conjecture fails?
\end{exercise}

\section{Hints}

\begin{hint}{conjectures:ex:linetwo}
Colour the half-open intervals $[k,k+1[$ alternately. For the plane, an equilateral triangle.
\end{hint}
\begin{hint}{conjectures:ex:spindle}
Three pairwise non-adjacent vertices would include at most one of $A,B,C$ and at most one of $A,E,F$. For the edges, use the symmetries of the spindle:
there are four kinds of edges.
\end{hint}
\begin{hint}{conjectures:ex:hexcolour}
$4(i^2+ij+j^2)=(2i+j)^2+3j^2$.
The neighbours are at $\pm a_1,\pm a_2,\pm(a_1-a_2)$. Two conditions are used: one on the diameter of a hexagon, one on the distance between
same-coloured hexagons.
\end{hint}
\begin{hint}{conjectures:ex:unitcube}
Two sums whose sets differ in one element are at distance $1$. Choose $v_1,v_2,\dots$ one at a time, avoiding finitely many bad directions.
\end{hint}
\begin{hint}{conjectures:ex:lenz}
$\|p-q\|^2=\|p\|^2+\|q\|^2$ when $p$ and $q$ are orthogonal.
\end{hint}
\begin{hint}{conjectures:ex:cupcap}
On a concave graph, every three points form a cap.
\end{hint}
\begin{hint}{conjectures:ex:pentwitness}
$X$ is inside the triangle $YZW$ when the three determinants $\det(Z-Y,X-Y)$, $\det(W-Z,X-Z)$, $\det(Y-W,X-W)$ have the same sign.
\end{hint}
\begin{hint}{conjectures:ex:ucsharp}
(c) Removing $\emptyset$ removes one set and no element. Compare $\frac aN$ and $\frac a{N-1}$.
\end{hint}
\begin{hint}{conjectures:ex:ucdual}
Take complements in $\bigcup\mathcal F$.
\end{hint}
\begin{hint}{conjectures:ex:ucgolden}
$\frac1{\varphi^2}=1-\frac1\varphi$, so $h(\varphi^{-2})=h(\varphi^{-1})$.
\end{hint}

\section{Solutions}

\begin{solution}{conjectures:ex:linetwo}
\idea A line has no odd cycles of unit steps that close up; a triangle does.
\writeup Colour $x\in\mathbb R$ by the parity of $\lfloor x\rfloor$. If $|x-y|=1$, then $\lfloor y\rfloor=\lfloor x\rfloor\pm1$, of the other parity. One
colour is not enough, since $0$ and $1$ are at distance $1$. In the plane, the three corners of an equilateral triangle of side $1$ are pairwise at
distance $1$, so they need three colours.
\end{solution}

\begin{solution}{conjectures:ex:spindle}
\idea The spindle is two triangles of triangles; each can host at most one vertex of an independent set, apart from $A$.
\writeup Let $I$ be a set of pairwise non-adjacent vertices. If $A\in I$, then $B,C,E,F\notin I$, and $D,G$ are adjacent, so $|I|\le2$. If $A\notin I$,
then $I$ has at most one of $B,C$ (adjacent) and at most one of $E,F$; if it has one of $B,C$ it cannot have $D$, and if it has one of $E,F$ it cannot
have $G$; and it cannot have both $D$ and $G$. So $I$ has at most one vertex from $\{B,C,D\}$ and at most one from $\{E,F,G\}$: $|I|\le2$. A
$3$-colouring splits the seven vertices into three such sets, of total size at most $6<7$.

The reflection exchanging the two rhombi and the reflections inside each rhombus (exchanging $B,C$ and $E,F$ together, or separately) show that there are
four kinds of edges: $DG$; the four edges at $A$; $BC$ and $EF$; the four edges $BD,CD,EG,FG$. Colour with $1,2,3$. Without $DG$: $A=1$, $B=E=2$,
$C=F=3$, $D=G=1$. Without $AB$: $A=1$, $B=1$, $C=2$, $D=3$, $E=2$, $F=3$, $G=1$; check $DG$: $3\ne1$, $EG$: $2\ne1$, $FG$: $3\ne1$. Without $BC$:
$A=1$, $B=C=2$, $D=3$, $E=2$, $F=3$, $G=1$. Without $BD$: $A=1$, $B=2$, $C=3$, $D=2$, $E=2$, $F=3$, $G=1$; check $CD$: $3\ne2$, $DG$: $2\ne1$. By
symmetry, every edge is covered.
\end{solution}

\begin{solution}{conjectures:ex:hexcolour}
\idea Complete the square; neighbouring colours differ by $\pm1,\pm3,\pm2$ modulo $7$.
\writeup If $i^2+ij+j^2=N$ then $(2i+j)^2+3j^2=4N$, so $3j^2\le4N$. For $N=2,5,6$, the numbers $4N-3j^2$ with $3j^2\le4N$ are $8,5$; $20,17,8$;
$24,21,12$, none a square. For $N=1$: $j^2\le1$; $j=0$ gives $2i=\pm2$, and $j=\pm1$ gives $(2i\pm1)^2=1$; this yields $\pm(1,0),\pm(0,1),\pm(1,-1)$. For
$N=3$: $(2i+j)^2+3j^2=12$ forces $j^2=1$, $(2i+j)^2=9$, or $j^2=4$, $2i+j=0$; this yields $\pm(1,1),\pm(2,-1),\pm(1,-2)$. For $N=4$: $16-3j^2$ is a square
for $j=0$ ($2i=\pm4$) and $j=\pm2$ ($2i+j=\pm2$), giving $\pm(2,0),\pm(0,2),\pm(2,-2)$; $j=\pm1$ gives $13$, not a square.
Moving to a neighbour changes $(i,j)$ by one of the six vectors $\pm(1,0)$, $\pm(0,1)$, $\pm(1,-1)$, so the colour changes by $\pm1$, $\pm3$
or $\pm2$ modulo $7$, six different nonzero residues: the hexagon and its neighbours use all seven colours. The proof needs $2s<1$ and $(\sqrt{21}-2)s>1$, so it works for
$\frac1{\sqrt{21}-2}<s<\frac12$, that is $0.3871\ldots<s<0.5$. (The true distance between same-coloured hexagons is $\sqrt7\,s$, a little more than the bound $(\sqrt{21}-2)s$
used here, as a direct computation gives; so the colouring works for $\frac1{\sqrt7}=0.378\ldots<s<\frac12$.)
\end{solution}

\begin{solution}{conjectures:ex:unitcube}
\idea The cube $\{0,1\}^k$ projected generically to the plane keeps its $k2^{k-1}$ edges of length $1$.
\writeup Two sums coincide when $\sum_i\varepsilon_iv_i=0$ with $\varepsilon_i\in\{-1,0,1\}$ not all $0$. Choose the $v_j=(\cos\theta_j,\sin\theta_j)$
one at a time. If $v_1,\dots,v_{j-1}$ have distinct subset sums, a bad relation involving $v_j$ has $\varepsilon_j=\pm1$ and fixes
$v_j=\mp\sum_{i<j}\varepsilon_iv_i$: finitely many forbidden points of the circle, so a good $\theta_j$ exists. Now the sums for $S$ and $S\cup\{i\}$
($i\notin S$) differ by $v_i$, of length $1$. There are $k2^{k-1}$ such pairs (choose $i$, then $S\subseteq\{1,\dots,k\}\setminus\{i\}$), that is
$\frac n2\log_2n$ for $n=2^k$.
\end{solution}

\begin{solution}{conjectures:ex:lenz}
\idea Two orthogonal circles of radius $\frac1{\sqrt2}$ are at distance $1$ from each other everywhere.
\writeup For $p=(x,y,0,0)$ and $q=(0,0,z,w)$ on the two circles, $\langle p,q\rangle=0$, so $\|p-q\|^2=\|p\|^2+\|q\|^2=\frac12+\frac12=1$. All
$\frac n2\cdot\frac n2$ pairs with one point on each circle are at distance $1$. So in $\mathbb R^4$ the unit distance problem has a quadratic answer, and
the plane's $n^{4/3}$ is a planar phenomenon.
\end{solution}

\begin{solution}{conjectures:ex:cupcap}
\idea Concave graphs only contain caps; the bound $g(3,l)\le l$ is the theorem.
\writeup By Theorem~\ref{conjectures:thm:cupcap}, $g(3,l)\le\binom{l-1}1+1=l$. Take $l-1$ points on the graph of $x\mapsto-x^2$ with distinct first coordinates.
The slope of the chord from $(s,-s^2)$ to $(t,-t^2)$ is $-(s+t)$, so for $s<t<r$ the slopes $-(s+t)>-(t+r)$ decrease: every three points form a $3$-cap and
none a $3$-cup, and there is no $l$-cap since there are only $l-1$ points. So $g(3,l)>l-1$. Reflecting in the horizontal axis exchanges cups and caps,
so $g(k,3)=g(3,k)=k$.
\end{solution}

\begin{solution}{conjectures:ex:pentwitness}
\idea Inside means on the same side of each edge.
\writeup Write $d(Y,Z,X)=\det(Z-Y,X-Y)$, the determinant of the two vectors as columns. For $Q_3=(2,-1)$ in $P_0P_1Q_2$, with $P_0=(6,0)$,
$P_1=(0,6)$, $Q_2=(-1,-2)$:
\[
\begin{aligned}d(P_0,P_1,Q_3)&=\det\bigl((-6,6),(-4,-1)\bigr)=6+24=30,\\ d(P_1,Q_2,Q_3)&=\det\bigl((-1,-8),(2,-7)\bigr)=7+16=23,\\
d(Q_2,P_0,Q_3)&=\det\bigl((7,2),(3,1)\bigr)=7-6=1,\end{aligned}
\]
all positive. For $Q_1=(-2,1)$ in $P_0P_2Q_0$, with $P_2=(-6,0)$ and $Q_0=(1,2)$:
\[
\begin{aligned}d(P_0,P_2,Q_1)&=\det\bigl((-12,0),(-8,1)\bigr)=-12,\\ d(P_2,Q_0,Q_1)&=\det\bigl((7,2),(4,1)\bigr)=7-8=-1,\\
d(Q_0,P_0,Q_1)&=\det\bigl((5,-2),(-3,-1)\bigr)=-5-6=-11,\end{aligned}
\]
all negative.
\end{solution}

\begin{solution}{conjectures:ex:ucsharp}
\idea Power sets are extremal; small sets are what make an element frequent.
\writeup (a) Complementation is a bijection between the subsets containing $x$ and those not containing it. (b) $\{\emptyset,\{1\},\{1,2\}\}$ is
union-closed and $2$ has frequency $\frac13$. $\{\{1\},\{2\},\{3\}\}$ is not union-closed, and each element has frequency $\frac13$. (c) The union of two nonempty sets is nonempty, so $\mathcal F\setminus\{\emptyset\}$ is union-closed. If $x$ lies in $a$ of the $N$ sets
of $\mathcal F$, it lies in $a$ of the $N-1$ sets of $\mathcal F\setminus\{\emptyset\}$, and $\frac a{N-1}>\frac aN$: frequencies go up, so the
conjecture for $\mathcal F$ gives it for $\mathcal F\setminus\{\emptyset\}$. Conversely, $a\ge\frac{N-1}2$ gives only
$\frac aN\ge\frac12-\frac1{2N}$. The empty set changes frequencies by less than $\frac1{2N}$, which is why some authors include it and some do not.
\end{solution}

\begin{solution}{conjectures:ex:ucdual}
\idea Complementation turns unions into intersections and frequency $p$ into $1-p$.
\writeup Let $X=\bigcup\mathcal F$ and $\mathcal G=\{X\setminus S:S\in\mathcal F\}$. Then $(X\setminus S)\cap(X\setminus T)=X\setminus(S\cup T)$, so
$\mathcal F$ is union-closed if and only if $\mathcal G$ is intersection-closed, and $x\in X$ has frequency $p$ in $\mathcal F$ exactly when it has
frequency $1-p$ in $\mathcal G$. If $\mathcal F$ has a single set, it is $X\neq\emptyset$ and the conjecture is trivial; otherwise $\mathcal G$ has
at least two sets. Conversely, an intersection-closed $\mathcal G$ with at least two sets has $X=\bigcup\mathcal G\ne\emptyset$, and taking complements
in $X$ gives back a union-closed $\mathcal F$ with at least two different sets, so with a nonempty one. So ``some $x$ has frequency at least $\frac12$
in $\mathcal F$'' becomes ``some $x\in X$ has frequency at most $\frac12$ in $\mathcal G$''.
\end{solution}

\begin{solution}{conjectures:ex:ucgolden}
\idea At the golden point, squaring a probability is the same as taking its complement.
\writeup (a) Let $x=\frac1\varphi$. Since $\varphi^2=\varphi+1$, $x^2=\frac1{\varphi^2}=1-\frac1\varphi=1-x$, so $h(x^2)=h(1-x)=h(x)$. And
$\varphi\,x\,h(x)=h(x)$. (b) For $U=V=x$ constant, the hypothesis reads $h(x^2)\ge c\,x\,h(x)$, that is $h(x)\ge\frac c\varphi h(x)$, false for $c>\varphi$
since $h(x)>0$. (c) The theorem only gives a sufficient condition; a better argument may use more than the one-coordinate inequality, for instance
the fact that the $U_i$ come from a union-closed family, which the hypothesis ignores. The remark's improvements beyond $\frac{3-\sqrt5}2$ do exactly this.
\end{solution}

\section*{Sources and further reading}

\begin{itemize}
\item A.~Soifer, \emph{The Mathematical Coloring Book}, Springer, 2009: the history of the Hadwiger--Nelson problem, including Nelson's question and
Isbell's colouring in 1950. M.~Gardner, \emph{Mathematical Games}, Scientific American, October 1960.
\item L.~Moser and W.~Moser, \emph{Solution to Problem 10}, Canadian Mathematical Bulletin 4 (1961), 187--189.
\item A.~D.~N.~J.~de Grey, \emph{The chromatic number of the plane is at least 5}, Geombinatorics 28 (2018), 18--31; arXiv:1804.02385, which
describes the $20\,425$-vertex graph and the $1581$-vertex one. E.~Lamb, \emph{Decades-old graph problem yields to amateur mathematician}, Quanta
Magazine, April 2018.
\item J.~Parts, \emph{Graph minimization, focusing on the example of 5-chromatic unit-distance graphs in the plane}, Geombinatorics 29 (2020),
137--166.
\item K.~J.~Falconer, \emph{The realization of distances in measurable subsets covering $\mathbb R^n$}, Journal of Combinatorial Theory A 31 (1981),
184--189.
\item P.~Erd\H os, \emph{On sets of distances of $n$ points}, American Mathematical Monthly 53 (1946), 248--250.
\item J.~Spencer, E.~Szemer\'edi and W.~T.~Trotter, \emph{Unit distances in the Euclidean plane}, in \emph{Graph Theory and Combinatorics}, Academic
Press, 1984, 293--303. L.~A.~Sz\'ekely, \emph{Crossing numbers and hard Erd\H os problems in discrete geometry}, Combinatorics, Probability and
Computing 6 (1997), 353--358.
\item P.~Erd\H os and G.~Szekeres, \emph{A combinatorial problem in geometry}, Compositio Mathematica 2 (1935), 463--470; and \emph{On some extremum
problems in elementary geometry}, Annales Univ.\ Sci.\ Budapest 3--4 (1960--61), 53--62.
\item W.~Morris and V.~Soltan, \emph{The Erd\H os--Szekeres problem on points in convex position: a survey}, Bulletin of the AMS 37 (2000), 437--458.
G.~Szekeres and L.~Peters, \emph{Computer solution to the 17-point Erd\H os--Szekeres problem}, ANZIAM Journal 48 (2006), 151--164.
A.~Suk, \emph{On the Erd\H os--Szekeres convex polygon problem}, Journal of the AMS 30 (2017), 1047--1053.
\item H.~Bruhn and O.~Schaudt, \emph{The journey of the union-closed sets conjecture}, Graphs and Combinatorics 31 (2015), 2043--2074;
arXiv:1309.3297: the origin of the conjecture and the small cases.
\item J.~Gilmer, \emph{A constant lower bound for the union-closed sets conjecture}, arXiv:\allowbreak2211.09055 (2022). R.~Alweiss, B.~Huang and M.~Sellke,
\emph{Improved lower bound for Frankl's union-closed sets conjecture}, arXiv:2211.11731. Z.~Chase and S.~Lovett, arXiv:2211.11689. W.~Sawin,
\emph{An improved lower bound for the union-closed set conjecture}, arXiv:2211.11504. J.~Liu, arXiv:2306.08824 (2023).
\item R.~B.~Boppana, \emph{A useful inequality for the binary entropy function}, arXiv:2301.09664 (2023).
\item Lemma~\ref{conjectures:lem:hxy}, with the constant $\frac1{1+\ln2}$ and the series proof, was written for this book as a short substitute for Boppana's
inequality; it gives the frequency bound $\frac{1-\ln2}2$, between Gilmer's $0.01$ and the golden-ratio constant.
\end{itemize}
